\exercisesection{光滑代数}

\begin{enumerate}
\def\labelenumi{\arabic{enumi})}
\tightlist
\item
  设 \(k\) 为环，A 为一个线性拓扑的形式光滑 \(k\)-代数（带线性拓扑）。若 A 满足下列条件，则称 A 是形式平展（相应地，形式非分歧）的 \(k\)-代数：无论 \(k\)-代数 C（赋予离散拓扑）与 C 的平方零理想 N 如何，每个 \(k\)-代数的连续同态 A → C/N 都容许唯一（相应地，至多一个）提升 A → C。为使 A 在 \(k\) 上形式平展，必须且只需它既形式光滑又形式非分歧。\(k\) 的每个商都是形式非分歧的。
\end{enumerate}

\begin{enumerate}
\def\labelenumi{\alph{enumi})}
\item
  对形式光滑或形式平展代数证明命题 3 与命题 4 的类似陈述。
\item
  设 J 为 \(\Lambda\) 的一个理想，使 A 的拓扑是 J-进拓扑；用 \(\widehat{\Omega}_{k}(\Lambda)\) 表示 A-模 \(\Omega_{k}(A)\)（关于 J-进拓扑）的分离完备化。为使 \(k\)-拓扑代数 A 形式光滑，必须且只需 \(\widehat{\Omega}_{k}(A)\) 为零（注意这两个性质都等价于：A 到每个被 J 的某个幂零化的 A-模中的 \(k\)-导子均为零）。
\end{enumerate}

\begin{enumerate}
\def\labelenumi{\arabic{enumi})}
\setcounter{enumi}{1}
\item
  设 A 为 Noether 环，J 为 A 的理想，\(\widehat{\mathbf{A}}\) 为 A 关于 J-进拓扑的分离完备化。若 \(\Lambda\)-代数 \(\widehat{\mathbf{A}}\)（关于离散拓扑）形式光滑，则它是形式平展的（注意有 \(\Omega_{\mathbf{A}}(\widehat{\mathbf{A}}) = J^{n}\Omega_{\mathbf{A}}(\widehat{\mathbf{A}})\) 对每个整数 \(n\) 成立）。
\item
  设 K 为域 \(k\) 的一个有限生成扩张；证明下列条件等价：
\end{enumerate}

\begin{enumerate}
\def\labelenumi{(\roman{enumi})}
\item
  k-代数 K 是形式光滑的；
\item
  \(k\)-代数 K 是形式平展的；
\item
  K 是 k 的有限可分扩张。
\end{enumerate}

\begin{enumerate}
\def\labelenumi{\arabic{enumi})}
\setcounter{enumi}{3}
\tightlist
\item
  设 \(k\) 为 Noether 环，A 为一个有限生成的 \(k\)-代数。若 \(\Lambda\) 赋予离散拓扑后是形式平展（相应地，形式光滑）的，则称它是 \(k\)-平展（相应地，光滑）代数。
\end{enumerate}

\begin{enumerate}
\def\labelenumi{\alph{enumi})}
\tightlist
\item
  证明下列条件等价：
\end{enumerate}

\begin{enumerate}
\def\labelenumi{(\roman{enumi})}
\item
  k-代数 A 是光滑的；
\item
  \(\Omega_k(A) = 0\)
\item
  \(\mathrm{A}\otimes_{k}\mathrm{A}\)-模 A 是投射的。
\end{enumerate}

（设 I 为典范同态 \(\mathrm{A} \otimes_k \mathrm{A} \to \mathrm{A}\) 的核；证明 (ii) 与 (iii) 都等价于

\begin{enumerate}
\def\labelenumi{(\roman{enumi})}
\setcounter{enumi}{3}
\tightlist
\item
  \(I_{q} = 0\) 对 A \(\otimes_{k}\) A 的每个包含 I 的素理想 q 成立。）
\end{enumerate}

\begin{enumerate}
\def\labelenumi{\alph{enumi})}
\setcounter{enumi}{1}
\item
  设 \(f \in k[T]\) 为一个使 f 与 \(f'\) 生成 k{[}T{]} 的单位理想的多项式。证明 k-代数 \(k[T]/(f)\) 是平展的（参见第 3 节，例 3）。
\item
  若 k 是域，则光滑 k-代数就是平展 k-代数，且此概念与 A, V, 第~28 页引入的概念一致（k 上有限次代数的情形由该处第~32 页定理 3 得到；若 A 是光滑 k-代数，考虑 A 的一个极大理想 m，并对 A/m \(^{n}\) 应用前一情形）。于是光滑 k-代数就是 k 的有限可分扩张的有限积（A, V, 第~34 页，定理 4）。
\item
  为使 k-代数 A 是光滑的，必须且只需对 k 的每个素理想 p，\(\kappa(\mathfrak{p})\)-代数 \(\kappa(\mathfrak{p})\otimes_{k}A\) 都是平展的。
\end{enumerate}

\begin{enumerate}
\def\labelenumi{\arabic{enumi})}
\setcounter{enumi}{4}
\tightlist
\item
  设 A 为完备 Noether 局部环，B 为一个本质有限型 A-代数。
\end{enumerate}

\begin{enumerate}
\def\labelenumi{\alph{enumi})}
\item
  证明 A 中使环 \(\Lambda_{\mathfrak{p}}\) 正则的素理想 p 所成的集合在 \(\operatorname{Spec}(\mathbf{A})\) 中是开的（将利用 VIII, § 5 习题 16 化为 A 是整环的情形，再利用 IX, § 2, 第 5 节, 定理 3 与 VIII, § 5 习题 17）。
\item
  假设 A 包含一个域 \(^{1}\)。证明 B 中使环 \(B_{q}\) 正则的素理想 q 所成的集合在 Spec(B) 中是开的（将利用 § 7, 第 9 节, 定理 3 的系 2 与 IX, § 3, 第 3 节, 定理 2）。
\end{enumerate}

\begin{enumerate}
\def\labelenumi{\arabic{enumi})}
\setcounter{enumi}{5}
\item
  设 \(k\) 为环，A 为一个线性拓扑的形式光滑 \(k\)-代数。证明对 A 的每个开理想 J，A/J-模 \((\mathrm{A} / \mathrm{J}) \otimes_{\mathrm{A}} \Omega_k(\mathrm{A})\) 是投射的（利用第 1 节的例子证明：对 A/J-模的每个满同态 \(\mathrm{M} \to \mathrm{M}''\)，每个 \(k\)-导子（自 A 到 \(\mathrm{M}''\)）都可提升到 M）。
\item
  设 \(k\) 为环，\(\rho : A \to B\) 为 \(k\)-代数同态，\(J\) 为 \(B\) 的理想。称 \(B\)（赋予 \(J\)-进拓扑）在 \(A\) 上关于 \(k\) 形式光滑，如果对每个 \(A\)-代数 \(C\)（赋予离散拓扑）与每个平方为零的理想 \(N\)（在 \(C\) 中），每个 \(A\)-代数的连续同态 \(B \to C/N\)，只要它可提升为一个 \(k\)-代数连续同态（从 \(B\) 到 \(C\) 中），就也可提升为一个 \(A\)-代数连续同态（从 \(B\) 到 \(C\) 中）。
\end{enumerate}

\begin{enumerate}
\def\labelenumi{\alph{enumi})}
\tightlist
\item
  证明下列条件等价：
\end{enumerate}

\begin{enumerate}
\def\labelenumi{(\roman{enumi})}
\item
  B 关于 J-进拓扑在 A 上关于 k 形式光滑。
\item
  对每个 \(k\)-导子 D（从 \(\Lambda\) 到一个被 J 的某个幂零化的 B-模 E 中），存在 \(k\)-导子 \(\tilde{D}: B \to E\) 使 \(\tilde{D} \circ \rho = D\)；
\item
  对每个整数 \(n \geqslant 0\)，典范映射 \(u_n: \Omega_k(A) \otimes_A (B/J^n) \longrightarrow \Omega_k(B) \otimes_B (B/J^n)\) 是单射且其像是直和项。
\end{enumerate}

\begin{enumerate}
\def\labelenumi{\alph{enumi})}
\setcounter{enumi}{1}
\tightlist
\item
  假设 B 关于 J-进拓扑在 k 上形式光滑。证明：为使 B 在 A 上形式光滑，必须且只需典范同态 \(u_{1}:\Omega_{k}(A)\otimes_{A}(B/J)\to\Omega_{k}(B)\otimes_{B}(B/J)\) 是单射且其像是直和项（利用 a) 与习题 6 证明 \(u_{1}\) 的一个收缩可提升为每个 \(u_{n}\) 的收缩）。
\end{enumerate}

¶ 8) 设 k 为域，T 为一个有限的未定元族。给 k-代数 \(k[[T]]\) 赋予离散拓扑。

\begin{enumerate}
\def\labelenumi{\alph{enumi})}
\item
  证明 \(k[[T]]\) 在 \(k[T]\) 上关于 k 形式光滑（习题 7）。
\item
  若 \(k[[\mathbf{T}]]\) 在 \(k\) 上形式光滑，证明特征 \(p\)（\(k\) 的）严格为正且 \(k\) 是 \(k^p\) 的有限扩张（由习题 2 推出 \(\Omega_{k|\mathbf{T}|}(k[[\mathbf{T}]])\) 为零，从而 \(\Omega_{k(\mathbf{T})}(k((\mathbf{T})))\) 为零；然后利用 A, V, 第~99 页命题 6 证明：\(k((\mathbf{T}))\) 中由系数生成 \(k^p\) 上有限生成扩张的幂级数组成的中介扩张等于 \(k((\mathbf{T})))\)。
\end{enumerate}

\begin{enumerate}
\def\labelenumi{\arabic{enumi})}
\setcounter{enumi}{8}
\tightlist
\item
  设 \(k\) 为特征 \(p > 0\) 的域，包含子域的一个递减有向族 \((k_{\lambda})_{\lambda \in \Lambda}\)，满足 \(\bigcap_{\lambda \in \Lambda} k_{\lambda} = k^{p}\)。
\end{enumerate}

\begin{enumerate}
\def\labelenumi{\alph{enumi})}
\item
  证明 \(\Omega(k)\) 到 \(\Omega_{k_{\lambda}}(k)\) 中的典范同态（\(\lambda\) 取遍 \(\Lambda\)）的核的交化为 (0)（考虑一个 \(p\)-基，取自 \(k\)）。
\item
  设 A 为正则局部 k-代数，赋予 \(m_{A}\)-进拓扑。证明：为使 A 在 k 上形式光滑，必须且只需 \(\Lambda\) 在 k 上关于 \(k_{\lambda}\) 形式光滑对每个 \(\lambda \in \Lambda\) 成立（利用 a)、命题 8 与习题 7, a)）。
\end{enumerate}

¶ 10) 设 A 为完备 Noether 局部环，K 为其分式域。设 p 为 A 的素理想；用 B 表示局部环 \(A_{p}\)，\(\widehat{B}\) 表示其完备化。

\begin{enumerate}
\def\labelenumi{\alph{enumi})}
\item
  假设 A 正则且 K 的特征为零。证明 K-代数 \(K \otimes_{B} \widehat{B}\) 是绝对正则的。
\item
  假设 A 正则且 K 的特征 p \textgreater{} 0，于是 A 与形式幂级数环 \(k[[T_{1}, \ldots, T_{n}]]\) 等同（IX, § 3, 第 3 节, 定理 2）。用 \((k_{\lambda})_{\lambda \subset \Lambda}\) 表示子域 \(k_{\lambda}\)（k 的子域，包含 \(k^{p}\)，且使 k 是 \(k_{\lambda}\) 的有限次扩张）所成的族。设 \(\lambda \in \Lambda\)；用 \(A_{\lambda}\) 表示环 \(k_{\lambda}[[T_{1}^{p}, \ldots, T_{n}^{p}]]\)，\(B_{\lambda}\) 表示 \(A_{\lambda}\) 在素理想 \(p \cap A_{\lambda}\) 处的局部环，\(K_{\lambda}\) 表示 \(\Lambda_{\lambda}\) 与 \(B_{\lambda}\) 的分式域。证明 A-代数 B（赋予离散拓扑）关于 \(B_{\lambda}\) 形式光滑（可先证明 \(\widehat{B}\) 同构于 \(B \otimes_{B_{\lambda}} \widehat{B_{\lambda}}\)）。
\item
  在 b) 的假设下，证明 \(K \otimes_{B} \widehat{B}\) 的每个局部环都形如 \((\widehat{B})_{\mathbf{r}}\)，其中 \(\mathfrak{r} \in \operatorname{Spec}(\widehat{B})\) 且 \(\mathfrak{r} \cap B = \{0\}\)，并且它在 K 上关于 \(K_{\lambda}\) 形式光滑。由此推出 K-代数 \(K \otimes_{B} \widehat{B}\) 是绝对正则的（证明 K 的子域族 \(\{K_{\lambda}\}_{\lambda \in \Lambda}\) 是递减有向的，且 \(\bigcap_{\lambda \in \Lambda} K_{\lambda} = k^{p}((T_{1}, \ldots, T_{n}))\)；然后应用习题 9）。
\item
  回到一般情形。证明 A 是 Grothendieck 环（§ 6, 习题 5）；若 \(\mathfrak{q} \in \operatorname{Spec}(B)\)，需要证明 \(\kappa(\mathfrak{q})\)-代数 \(\kappa(\mathfrak{q}) \otimes_{\mathrm{B}} \widehat{\mathrm{B}}\) 是绝对正则的；把 A 换成 \(\mathrm{A}/(\mathrm{A} \cap \mathfrak{q})\) 后可假设 \(\Lambda\) 是整环且 \(\mathfrak{q} = \{0\}\)；将利用 IX, § 2, 第 3 节, 定理 3 与 § 3, 第 3 节, 定理 2 化为 A 是完备正则局部环的情形。
\end{enumerate}

\begin{enumerate}
\def\labelenumi{\arabic{enumi})}
\setcounter{enumi}{10}
\item
  设 A 为 Noether 环；假设对 A 的每个极大理想 m，局部环 \(A_{m}\) 的完备化都是绝对正则的 \(A_{m}\)-代数。证明 A 是 Grothendieck 环（§ 6, 习题 5）。（利用该处与前一习题证明：对 A 的每个极大理想 m，环 \(A_{m}\) 都是 Grothendieck 环。）
\item
  设 A 为 Noether 局部环。为使 A 是 Grothendieck 环，必须且只需：对每个有限整 A-代数 B、B 的每个极大理想 m 以及完备化环 \(\widehat{B_{m}}\) 的每个满足 \(B_{m} \cap q = (0)\) 的素理想 q，局部环 \((\widehat{B_{m}})_{q}\) 都是正则的（为证此条件必要，由前一习题只需证明：对 A 的每个素理想 p 与 \(\kappa(\mathfrak{p})\) 的每个有限扩张 K，环 \(K \otimes_{A} \widehat{A}\) 都是正则的；可引入一个以 \(K\) 为分式域的有限整 A-代数 B，记 \(\mathfrak{m}_{1},\ldots,\mathfrak{m}_{r}\) 为其极大理想，并把 \(K\otimes_{A}\widehat{A}\) 的局部环解释为某个 \(\widehat{B_{\mathfrak{m}_{i}}}\) 的局部环。
\item
  设 \(k\) 为域，A 为一个本质有限型 \(k\)-代数。证明 A 是 Grothendieck 环（化为 A 是 \(k\) 上有限个未定元的多项式环的局部环的情形；利用前一习题，再化为证明：对每个素理想 \(\mathfrak{r}\)（属于 \(\mathrm{A}[\mathrm{T}_1,\ldots ,\mathrm{T}_n]\)）、\(\mathrm{A}[\mathrm{T}_1,\ldots ,\mathrm{T}_n]\) 在包含 \(\mathfrak{r}\) 的极大理想处的每个局部环 C，以及每个理想 \(\mathfrak{q}\)（属于完备化环 \(\widehat{\mathrm{C}}\)，满足 \(\mathrm{C}\cap \mathfrak{q} = \mathfrak{r}\mathrm{C}\)），环 \((\widehat{\mathrm{C}})_q / \mathfrak{r}(\widehat{\mathrm{C}})_q\) 都是正则的。为此可利用第 9 节定理 3）。
\end{enumerate}

¶ 14) 设 A 为局部 Grothendieck 环。证明 A 中使环 \(A_{p}\) 正则的素理想 p 所成的集合在 \(\operatorname{Spec}(A)\) 中是开的（设 \(\widehat{A}\) 为 A 的完备化；将利用习题 5、VIII, § 5 习题 16，以及典范映射 \(\operatorname{Spec}(\widehat{A}) \to \operatorname{Spec}(A)\) 是开的这一事实）。

¶ 15) 设 \(k\) 为特征 \(p > 0\) 的完全域，\((\mathrm{T}_{n})_{n \in \mathbf{N}}\) 为一个未定元族，K 为域 \(k((\mathrm{T}_{n})_{n \in \mathbf{N}})\)。设 \(\Lambda\) 为形式幂级数环 \(\mathrm{K}[[\mathrm{X},\mathrm{Y}]]\) 的由 \(\mathrm{K}^{p}[[\mathrm{X},\mathrm{Y}]]\) 与 K 生成的子环。环 A 是维数 2 的 Noether 正则局部环，其完备化为 \(\mathrm{K}[[\mathrm{X},\mathrm{Y}]]\)。对每个整数 \(n\)，用 \(p_{n}\) 表示 A 的元素 \(\mathrm{X} - \mathrm{T}_{n}\mathrm{Y}\)，\(q_{n}\) 表示乘积 \(p_{0} \ldots p_{n}\)，并令 \(c = \sum_{n \in \mathbf{N}} q_{n}\mathrm{T}_{n}\)。用 B 表示 \(\mathrm{K}[[\mathrm{X},\mathrm{Y}]]\) 的由 \(c\) 生成的子 A-代数。

\begin{enumerate}
\def\labelenumi{\alph{enumi})}
\item
  设 \(n\) 为整数 \(> 0\)，\(\mathfrak{p}\) 为 \(B\) 的一个包含 \(p_n\) 的高度 1 的素理想。证明环 \(B_{\mathfrak{p}}\) 不是正规的（对每个整数 \(m\)，令 \(c_m = (c - \sum_{i=0}^{m-1} q_i T_i)/q_m\)，将证明 \(B_{\mathfrak{p}} = A_{A p_n}[c_{n-1}]\)，于是 \(c_n\) 不属于 \(B_{\mathfrak{p}}\)，而 \(c_n^p\) 属于 \(A\)）。
\item
  由此推出：B 中使环 \(B_{p}\) 正则的素理想 p 所成的集合不包含 Spec(B) 的任何非空开子集（注意 B 是忠实平坦 A-模）。
\end{enumerate}

